\begin{pdfdiagram}
  % 半对数绘图区；曲线端点与坐标轴箭头之间保留余量。
  \fill[BookCodeBg] (0,0) rectangle (8.25,4.35);
  \draw[BookRule, line width=0.45pt] (0,0) rectangle (8.25,4.35);
  \draw[BookMuted, line width=0.65pt, ->] (0,0) -- (8.85,0);
  \draw[BookMuted, line width=0.65pt, ->] (0,0) -- (0,4.75);
  \node[hw label, anchor=north] at (4.1,-0.62) {经过年数};
  \node[hw label, anchor=south west] at (0,4.80) {相对倍数（对数轴）};
  \foreach \x/\t in {0/0, 2.0/5, 4.0/10, 6.0/15, 8.0/20} {
    \draw[BookMuted, line width=0.5pt] (\x,0.07) -- (\x,-0.07);
    \node[hw label, anchor=north] at (\x,-0.11) {\t};
  }
  \foreach \y/\t in {0/1, 1.05/10, 2.10/100, 3.15/1000} {
    \draw[BookMuted, line width=0.5pt] (0.07,\y) -- (-0.07,\y);
    \node[hw label, anchor=east] at (-0.13,\y) {\t};
  }
  \foreach \x in {2,4,6,8} {
    \draw[BookRule, line width=0.25pt, dashed] (\x,0) -- (\x,4.35);
  }
  \foreach \y in {1.05,2.10,3.15,4.20} {
    \draw[BookRule, line width=0.25pt, dashed] (0,\y) -- (8.25,\y);
  }
  % 同一年的两条历史趋势。scope clip 是坐标图的硬边界。
  \begin{scope}
    \clip (0,0) rectangle (8.25,4.35);
    \fill[BookAccent!8] (0,0) -- (8.0,4.22) -- (8.0,0.71) -- cycle;
    \draw[BookAccent, line width=1.05pt] (0,0) -- (8.0,4.22);
    \draw[BookTeal, line width=1.05pt] (0,0) -- (8.0,0.71);
  \end{scope}
  \foreach \x/\ya/\yb in {2/1.055/0.178,4/2.11/0.355,6/3.165/0.533,8/4.22/0.71} {
    \fill[BookAccent] (\x,\ya) circle (1.5pt);
    \fill[BookTeal] (\x,\yb) circle (1.5pt);
  }
  \node[hw label, anchor=west, text=BookAccent] at (8.12,4.22) {算力 $\times3325$};
  \node[hw label, anchor=west, text=BookTeal] at (8.12,0.71) {带宽 $\times3.87$};
  \draw[{Stealth[length=1.8mm]}-{Stealth[length=1.8mm]}, BookRed, line width=0.65pt]
    (7.72,0.71) -- (7.72,4.22);
  \node[hw label, fill=white, inner sep=2pt] at (6.55,2.47)
    {同一年竖直距离\\$=$ 算力/带宽差距};
  \node[hw label, anchor=west] at (5.1,1.10) {20 年：约 859 倍};
  \node[hw label, anchor=west] at (0.45,0.43) {cache：缩短可见时延};
  \node[hw label, anchor=west] at (2.7,1.23) {预取、MSHR：增加在途请求};
  \node[hw label, anchor=west] at (5.05,2.03) {多核：合并在途能力};
\end{pdfdiagram}
